
\section{Introduction}

Large language models (LLMs) are an emerging class of machine learning (ML) systems geared toward learning from huge text-based datasets.
LLMs such as GPT-3\cite{brownLanguageModelsAre2020} have revolutionized the ability of Artificial Integillence (AI) systems to understand and generate human language.
Due to innovative changes in model architecture, training methods, and through the help of the popularity of online services like ChatGPT~\cite{rayChatGPTComprehensiveReview2023}, the field of LLMs is evolving rapidly.

The number of everyday users is also growing rapidly due to the myriad of use cases from translation~\cite{yaoBenchmarkingLLMbasedMachine2024}, classification~\cite{sunTextClassificationLarge2023}, code generation~\cite{liuYourCodeGenerated2023} to healthcare~\cite{clusmannFutureLandscapeLarge2023}.
Additionally, cloud-based LLM services are currently the go-to method of access to LLMs for everyday users, but as the availability of open-source LLMs and datasets increases, especially over the last few years, the need for edge-based, localized access and execution of LLMs has become more sought after.
Massive community-driven pushes have facilitated easy access to LLMs and rapid prototyping of new models and optimizations to enable efficient LLM inference on edge devices.
At the forefront of these pushes is the GPT-Generated Model Language~\cite{GGML} (GGML).
GGML is a tensor library for ML specialized in enabling large models and high performance on commodity hardware.
Furthermore, the \GGML's \LPP~project~\cite{LLP} is specialized towards running LLMs on edge devices, supporting LLM inference on commodity CPUs and GPUs.

Unfortunately, LLMs can be very computationally demanding, even for inference.
In addition, due to their large memory footprint, they require high memory capacity and bandwidth.
These properties of LLMs make them challenging to execute on resource-constrained edge devices.
For example, running LLMs on mobile phones or Internet-of-Things devices (IoT) devices is in some cases impossible due to memory constraints.
Hence, there is a great demand for developing and deploying custom hardware accelerators to efficiently run these LLMs on resource-constrained edge devices.
Fortunately, Field Programmable Gate Arrays (FPGAs) are ideal for designing new flexible and power-efficient accelerators that can take advantage of LLM optimizations, such as block floating point quantization.
While some FPGA-based accelerators~\cite{Khan2021NPE,Lu2020Transformer} already exist for LLM inference at the edge, with constant changes to LLM architectures and optimizations we are in need of new specialized FPGA-based accelerators.

To create new and innovative FPGA-based accelerator architectures for LLM inference at the edge, we need ways to quickly prototype and evaluate LLM-based inference accelerators to reduce development costs and increase design space exploration.
Hence, we propose SECDA-LLM, a new platform for designing, integrating and deploying specialized accelerators for LLMs at the edge.
SECDA-LLM employs the SECDA design methodology~\cite{harisSECDAEfficientHardware2021}, and similar to SECDA-TFLite~\cite{harisSECDATFLiteToolkitEfficient2023a}, it provides the user with the ability to quickly prototype accelerator designs with the target application framework, in this case, \LPP~project.
% Our SECDA-LLM platform enables the designer to consider hardware-software co-design optimizations~\cite{gibson_dlas_2024} from different levels of LLM execution and makes deployment of LLMs through FPGA-based accelerators effortless.
Our SECDA-LLM platform enables the designer to consider hardware-software co-design optimizations in terms of both algorithmic and hardware implementations~\cite{gibson_dlas_2024}, and makes deployment of LLMs through FPGA-based accelerators effortless.
The contributions of this work are as follows:

\begin{itemize}
    \item \textit{SECDA-LLM}, a new design platform using the SECDA methodology which enables the design, integration and deployment of FPGA-based accelerators for LLMs on resource-constrained edge devices.

    \item A case study to demonstrate \textit{SECDA-LLM}, where we prototype and deploy a new accelerator to efficiently execute quantized MatMul operations during LLM inference.
    %with the focus on model-specific optimizations.

    \item Evaluation of our initial accelerator design executing the TinyLlama model~\cite{zhangTinyLlamaOpenSourceSmall2024} on the PYNQ-Z1~\cite{pynqz1} board, where we achieve a $11\times$ speedup over dual-core ARM NEON-based CPU execution.
    
    % \item Evaluation of our initial accelerator design, where we evaluated the TinyLlama model~\cite{zhangTinyLlamaOpenSourceSmall2024} using our accelerator, achieving the initial speed of 20 seconds per token.
        
    % \item Evaluation of our initial accelerator design, where we achieve a speed of up to \uptooverallX and a \energyX energy reduction over-optimized ARM-NEON-based inference.    
\end{itemize}




% While custom accelerators already exist for Deep Neural Network (DNN) inference on the edge, with LLMs' scale and algorithmic differences, we need new specialized FPGA-based accelerators for them.
% While custom accelerators already exist for Deep Neural Network (DNN) inference on the edge, with LLMs' scale and algorithmic differences, we need new specialized FPGA-based accelerators for them.

% This demand already exists for Deep Neural Network (DNN) inference on the edge, but with LLMs' scale and algorithmic differences, we need new specialized accelerators for them.

% Fortunately, FPGAs are ideal for designing new flexible and power-efficient accelerators that can take advantage of LLM optimizations, such as quantization, for edge inference.